
\section{Sentence Objects}\fbeg{}\begin{frame}{\ft{Sentence Objects}}

\begin{sftxt}[.94\textwidth]

There are numerous XML-based and other markup formats for text 
documents, including JATS, BioC, TEI, S2ORC, and TAGML.  
Any markup format, however, will encode natural-language 
text according to generic structures involving 
trees or graphs, rather than the inherent structure 
of written text itself.  For instance, XML trees 
are defined in terms of parent and child nodes, siblings, 
elements' text content, and attribute values.  
Although there are analogous structures in natural 
language --- for instances, sentences in sequence 
are similar to sibling XML nodes many textual 
phenomena can be modeled only partially via 
XML structure alone.\\\vspace*{-13pt}

For these reasons, any markup file should be seen as 
the \textit{serialization} of textual content, not as 
the definitive structure of that content.  Conceptually, 
text elements are best understood as typed 
values in an Object-Oriented context.  Moreover, I believe 
sentence objects --- that scale of representation --- 
should be the foundation of such object systems.  Sentence 
objects are also contained inside larger elements 
(paragraphs, sections, subsections) and can encompass 
smaller parts (keywords, annotated fragments), but 
overall \makebox{text-object systems should pivot 
around the sentence level.}
  

\end{ftxt}

\end{frame}

